\begin{center}
	\textbf{ABSTRAK}
\end{center}

\begin{center}
	Integrasi \textit{Explainable AI} (XAI) pada Model \textit{Convolutional Neural Network}\\
	untuk Deteksi Tumor Otak Menggunakan Citra CT \textit{Scan}
	
	\vspace{0.3cm}
	oleh
	
	\vspace{0.3cm}
	Micky Valentino\\
	18222093
\end{center}

\vspace{0.5cm}

Tumor otak memiliki urgensi diagnosis yang tinggi, namun ketersediaan \textit{Magnetic Resonance Imaging} (MRI) di Indonesia masih terbatas, sehingga citra \textit{Computed Tomography} (CT) \textit{Scan} menjadi modalitas alternatif utama. Model \textit{deep learning} untuk diagnosis umumnya bersifat \textit{black box}, dan belum jelas sejauh mana arsitektur \textit{Convolutional Neural Network} (CNN) yang berbeda memengaruhi kualitas penjelasan visual, bukan hanya akurasi klasifikasinya. Tugas akhir ini mengembangkan dan membandingkan tiga arsitektur, yaitu \textit{Correlation Learning Mechanism} (CLM), \textit{Explainable CNN}, dan EfficientNetB0, untuk mendeteksi tumor otak pada 557 citra CT \textit{Scan} (393 tumor, 164 tanpa tumor). Model dengan performa terbaik diintegrasikan dengan \textit{Grad-CAM}, yang kesesuaiannya terhadap \textit{mask} tumor dievaluasi menggunakan metrik XAlign. CLM dan \textit{Explainable CNN} berulang kali mengalami kolaps kelas akibat keterbatasan data. EfficientNetB0 mencatat akurasi tertinggi pada satu skema pengujian (96,63\%), tetapi turun menjadi 85,21\% $\pm$ 11,19\% saat divalidasi dengan \textit{10-fold cross-validation}, menandakan performa yang belum stabil. Skor XAlign pada \textit{Grad-CAM} model tersebut hanya 0,1749, menunjukkan \textit{heatmap} yang menyebar jauh melampaui batas tumor sebenarnya. Temuan ini mengindikasikan adanya \textit{trade-off} antara akurasi dan kualitas penjelasan, sehingga hasil deteksi belum dapat divalidasi secara klinis hanya berdasarkan \textit{Grad-CAM}.

\vspace{0.5cm}
\noindent\textbf{Kata kunci:} \textit{Convolutional Neural Network}, \textit{Explainable AI}, \textit{Grad-CAM}, Tumor Otak, CT \textit{Scan}.